\begin{abstract}
We prove that an unmarked, ungraded finite local algebra associated
with a complex quadratic pencil determines the pencil up to congruence.
The truncation order $n^2+2n-4$ is sharp, and no regularity assumption
is imposed. Multiplication recovers the first relation, its determinant
rulings, and the orientation that identifies the pencil coefficient
line. The construction extends to families over arbitrary complex
bases. On the effectively rigidified failure stack, the recovered
source carries a scalar-normalized apolar algebra bundle. The associated
multiple-incidence boundary admits explicit multi-Rees equations and
relative lifting conditions. The sharp inverse is proved independently
of this boundary application; its relative, covering, recognition,
and spectral extensions are retained in the appendices. The paired boundary theory identifies the full
multiple-reduced-incidence modification and computes arbitrary-order
Jordan contact slices, thick Hilbert tails, transverse graph
singularities and covered stable-map tails. This master preserves
both mathematical bodies and is not a third submission.
\end{abstract}
\maketitle
\section{Introduction}
\label{sec:introduction-v153}
The first paper proves the sharp unmarked inverse for every complex
quadratic pencil. The second identifies the full normalized Hilbert
modification over the reduced corank-two incidence open and computes
nonreduced-contact slices of every order. The respective principal
results are Theorems~\ref{thm:artin-local-inverse-v146},
\ref{thm:multiple-incidence-v161}, and
\ref{thm:contact-hilbert-v161}. The singularity and stable-map
comparison are Proposition~\ref{prop:contact-deformations-v161}
and Theorem~\ref{thm:stable-tail-v161}. Their effective stack and
relative lifting interpretation is given in
Theorem~\ref{thm:boundary-stack-v161} and
Proposition~\ref{prop:relative-obstructions-v161}.

This master contains all earlier statements and proofs. The focused
papers place their core arguments before supplementary extensions;
no extension is silently used as a premise of the sharp inverse.
Classical arrangement blow-ups, complete quadrics, hypersurface
singularities and symmetric canonical forms are credited at their
points of use. The source-normalized envelope is not a subquotient
of the finite failure algebra, and a closed algebra does not choose
a normal deformation direction. The Ballico 1993 comparison remains
documentary-limited \cite{Ballico93}; no priority inference is drawn
from unavailable text.

